\documentclass{article}
\usepackage[utf8]{inputenc}
\usepackage[russian]{babel}
\usepackage{amsmath}
\usepackage{amssymb}
\usepackage{mathtools}
\usepackage[a4paper, left=1.5cm, right=1.5cm, top=1.5cm, bottom=1.5cm]{geometry}\setlength{\parindent}{0pt}

\title{ Билет 9. Дифференцируемость и интегрируемость интегралов, зависящих от параметра.}
\author{}
\date{}

\begin{document}

\maketitle

\section*{Дифференцируемость интегралов, зависящих от параметра}

Пусть $f(x,y)$ определена на $[a,b] \times [c,d]$.

Пусть $\forall y \in [a,b] \quad \exists \int_a^b f(x,\bar{y})\,dx$

Пусть $\exists f_y'(x,y)$ — непрерывна как функция двух переменных.

Тогда:
\[
G(y) = \int_a^b f(x,y)\,dx
\]

дифференцируема на $[c,d]$ и

\[
\exists G'(y) = \int_a^b f_y'(x,y)\,dx
\]

\subsection*{Доказательство:}

Фиксируем $\forall y_0 \in [c,d]$. $\exists G'(y_0)$.

По определению производной:
\[
\frac{G(y_0 + h) - G(y_0)}{h}, \quad h \ne 0, \quad y_0 + h \in [c,d]
\]

Так как $\exists f_y'(x,y)$, то по формуле конечных приращений Лагранжа:

\[
\frac{G(y_0 + h) - G(y_0)}{h}
= \int_a^b \frac{f(x,y_0 + h) - f(x,y_0)}{h}\,dx 
= \int_a^b f_y'\bigl(x,\, y_0 + \theta h\bigr)\,dx,
\quad \text{где } 0 < \theta < 1
\]
(Замечание: $\theta$ может зависеть от $x$ и $h$, но мы не будем этого выписывать явно — достаточно того, что $0 < \theta < 1$.)

Так как \quad $f_y'(x, y_0 + \theta h) \mathop{\rightrightarrows}\limits^{[a, b]}_{\,h \to 0\,} f_y'(x, y_0)$ (по непрерывности $f_y'$),  
и так как $f_y'(x,y)$ непрерывна на $[a,b] \times [c,d]$, то она \textit{равномерно непрерывна}:

\[
\forall \varepsilon > 0\ \exists \delta > 0:\ 
\forall y \in Y:\ |y - y_0| < \delta \quad \Rightarrow \quad \forall x \in [a,b]:\ 
\bigl| f_y'(x, y) - f_y'(x, y_0) \bigr| < \varepsilon.
\]

А это условие того, что 
\[
f(x,y) \xRightarrow[y \to y_0]{[a,b]} f(x,y_0).
\]

Нам нужно показать, что 
\[
f_y'\bigl(x,\,y_0 + \theta h\bigr) \xRightarrow[y \to y_0]{[a,b]} f_y'(x, y_0),
\]
а это значит, что
\[
\forall\, \varepsilon > 0\ \exists\, \delta > 0:\ 
\forall\, h \text : 0 < |h| < \delta,\ 
\forall\, x \in [a,b]\colon\ 
\bigl| f_y'(x, y_0 + \theta h) - f_y'(x, y_0) \bigr| < \varepsilon.
\]

Возьмём в качестве этого $\delta$ то $\delta$, которое гарантируется определением равномерной непрерывности $f_y'$ на $[a,b] \times [c,d]$.  
Поскольку $|h| < \delta$ и $0 < \theta < 1$, то
\[
|y_0 + \theta h - y_0| = |\theta h| < |h| < \delta,
\]
а значит,
\[
f_y'\bigl(x,\,y_0 + \theta h\bigr) \xRightarrow[h \to 0]{[a,b]} f_y'(x, y_0).
\]

Следовательно, по теореме о предельном переходе под знаком интеграла (для равномерной сходимости на конечном отрезке),
\[
\lim_{h \to 0} \int_a^b f_y'\bigl(x,\, y_0 + \theta h\bigr)\,dx 
= \int_a^b \lim_{h \to 0} f_y'\bigl(x,\, y_0 + \theta h\bigr)\,dx 
= \int_a^b f_y'(x, y_0)\,dx,
\]
то есть
\[
G'(y_0) = \int_a^b f_y'(x, y_0)\,dx.
\]

Теорема доказана.

\subsection*{Пример:}

\[
\int_1^2 \ln(x^2 + y^2)\,dx, \quad x \in [1,2], \quad y \in [0,1]
\]

(функция непрерывна)

\[
f_y'(x,y) = \frac{2y}{x^2 + y^2}
\]

\[
G'(y) = \int_1^2 \frac{2y}{x^2 + y^2}\,dx = \int_1^2 \frac{2\,d\left(\frac{x}{y}\right)}{\left(\frac{x}{y}\right)^2 + 1} =
\]

\[
= 2 \arctan\left(\frac{x}{y}\right) \Big|_1^2 = 2 \arctan\frac{2}{y} - 2 \arctan\frac{1}{y}
\]

\section*{Интегрируемость интегралов, зависящих от параметра}

\subsection*{Теорема об интегрировании интеграла по параметру}

Пусть $f(x,y)$ определена на $[a,b] \times [c,d]$ и непрерывна как функция двух переменных.

Положим:
\[
G(y) = \int_a^b f(x,y)\,dx
\]

Тогда:
\[
\exists \int_c^d G(y)\,dy = \int_c^d \int_a^b f(x,y)\,dx\,dy = \int_a^b \int_c^d f(x,y)\,dy\,dx
\]

\subsection*{Доказательство:}

Так как $f(x,y)$ непрерывна как фнкция двух переменных, то она интегрируема по любому прямоугольнику.

Так как $f$ непрерывна по параметру $y$, то $G(y)$ непрерывна на $[c,d]$ (по предыдущей теореме, если бы мы знали дифференцируемость, но здесь достаточно непрерывности).

Так как $G(y)$ непрерывна, то $\exists \int_c^d G(y)\,dy$.

Аналогично, $\exists \int_a^b \int_c^d f(x,y)\,dy\,dx$.

Нужно доказать равенство:

Рассмотрим общую задачу:
\[
F_1(z) = \int_c^z \int_a^b f(x,y)\,dx\,dy, \quad F_2(z) = \int_a^b \int_c^z f(x,y)\,dy\,dx
\]

Покажем, что $F_1'(z) = F_2'(z)$.

По теореме о дифференцировании интеграла по параметру:

\[
F_1'(z) = \int_a^b f(x,z)\,dx
\]

\[
F_2'(z) = \int_a^b \frac{\partial}{\partial z} \left( \int_c^z f(x,y)\,dy \right) dx = \int_a^b f(x,z)\,dx
\]

Следовательно, $F_1'(z) = F_2'(z)$ для всех $z \in [c,d]$.

Значит, $F_1(z) = F_2(z) + A$ для некоторой константы $A$.

Но $F_1(c) = 0$, $F_2(c) = 0$, следовательно, $A = 0$.

Таким образом, $F_1(z) = F_2(z)$ для всех $z \in [c,d]$.

В частности, при $z = d$ получаем требуемое равенство.

Теорема доказана.

\end{document}